\documentclass[main.tex]{subfiles}


\begin{document}

%from pagenumber to up
% \phantomsection 
% \hypertarget{toc}{}

 

 

 
   

\section{Дифракция света}


\textbf{Дифракция света}~--- это огибание световыми волнами краев препятствия. (Дифракция свойственна волнам любой природы.)

Пусть волна падает на экран с достаточно узкой щелью (рис.\,\ref{fig:p192}).

    \begin{figure}[H]
    \centering

    \begin{tikzpicture}[font=\small,            line cap=round,                             >={Stealth[inset=0pt,round,length=3mm, angle'=20]},vec/.style={>={Stealth[inset=0pt,round,length=3.5mm, angle'=20]}}]


%line
\def\h{1.6}
\foreach \i in{0,...,5}
\draw[red,xshift={-.2*\i cm}] (-.2,-\h) --++ (0,{2*\h});

\foreach \i in{1,...,10}
\draw[red] (0,0) ++ (-30:{\i*.2}) arc [radius={\i*.2}, start angle=-30, end angle=30];


%plane
\draw[thick] (0,0.1) -- (0,\h) node[below right,inner ysep=0] {Э};
\draw[thick] (0,-0.1) -- (0,-\h);


%vec
\draw[->,red] (-1.4,1.25) ++ (-1,0) --++ (1,0);
\draw[->,red] (-1.4,-1.25) ++ (-1,0) --++ (1,0);


    \end{tikzpicture}
\caption{Дифракция на щели}
\label{fig:p192}
\end{figure}

В данном случае на экран~Э слева падает \emph{плоская} волна: ее <<гребни>> выстроены вдоль параллельных линий (обозначенных красным цветом). После прохождения щели волна становится практически \emph{сферической}: <<гребни>> лежат на дугах концентрических окружностей (волна расходится)! Расходящаяся волна заходит в область предположительной тени от экрана. 

% Особый интерес представляет случай падения плоской световой волны на \textbf{дифракционную решетку}~--- совокупность очень узких щелей, разделенных непрозрачными промежутками (рис.\,\ref{fig:p193}). 
Пусть на \textbf{дифракционную решетку}~--- совокупность узких щелей, разделенных непрозрачными промежутками,~--- падает плоская волна света (рис.\,\ref{fig:p193}). 

\tikzfading[name=fade out,
inner color=transparent!50,
outer color=transparent!100]

    \begin{figure}[H]
    \centering

    \begin{tikzpicture}[font=\small,            line cap=round,                             >={Stealth[inset=0pt,round,length=3mm, angle'=20]},vec/.style={>={Stealth[inset=0pt,round,length=3.5mm, angle'=20]}}]


%grat
\foreach \i in{0,...,3}
\draw[thick,yshift={.7*\i cm}] (0,0.1) --++ (0,.5) coordinate (a);
\node[below right,inner ysep=0] at (a) {Р};

\foreach \i in{0,...,3}
\draw[thick,yshift={-.7*\i cm}] (0,-0.1) --++ (0,-.5) coordinate (b);


%d
\draw[densely dashed] (b) --++ (.5,0) coordinate (db);
\draw[densely dashed] ([yshift=.7 cm]b) --++ (.5,0) coordinate (da);
\draw[densely dashed] ([xshift=-.2 cm]da) --++ (0,-.7) node[midway,right] {$d$};


%line
\foreach \i in{1,...,5}
\draw[red] ([xshift={-.2*\i cm}]a) -- ([xshift={-.2*\i cm}]b);


%plane
\def\xs{4}
\draw[thick,gray] ([xshift=\xs cm]a) coordinate (aa) node[below right,inner ysep=0,black] {Э} -- ([xshift=\xs cm]b) coordinate (bb);


%rays path
\path (a);
\pgfgetlastxy{\ix}{\iy}
\foreach \i in{-2,-1,...,2}
\path[->,red] (0,0) coordinate (ce) --++ ({\xs},{\iy/3*(\i)})
coordinate (m\i)
;


%m
\node[right] at (m-2) {2};
\node[right] at (m2) {2};

\node[right] at (m-1) {1};
\node[right] at (m1) {1};

\node[right] at (m0) {$O$};


%shade
\begin{scope}
    \def\rs{.4}
    \clip (a) -- (aa) -- (bb) -- (b) -- cycle;
    \foreach \i in{-2,-1,...,2}
    \fill [red,path fading=fade out] (m\i) ++ ({-(\rs-abs(.1*\i))},{-(\rs-abs(.1*\i))}) rectangle++ ({2*(\rs-abs(.1*\i))},{2*(\rs-abs(.1*\i))});
\end{scope}


%rays
\path (a);
\pgfgetlastxy{\ix}{\iy}
\foreach \i in{-2,-1,...,2}
\draw[->,red] (0,0) coordinate (ce) --++ ({\xs},{\iy/3*(\i)})
coordinate (m\i)
;


%angles
\pic [draw,angle radius=15mm,angle eccentricity=1.2,"$\varphi_1$"] {angle = m0--ce--m1};
\pic [draw,angle radius=25mm,angle eccentricity=1.15] {angle = m0--ce--m2};
\pic [draw,angle radius=26mm,angle eccentricity=1.15,] {angle = m0--ce--m2};
\coordinate (mm) at (\xs,1.1);
\pic [angle radius=26mm,angle eccentricity=1.12,"$\varphi_2$"] {angle = m0--ce--mm};

\pic [draw,angle radius=16mm,angle eccentricity=1.19,"$\varphi_1$"] {angle = m-1--ce--m0};
\pic [draw,angle radius=27mm,angle eccentricity=1.15] {angle = m-2--ce--m0};
\pic [draw,angle radius=28mm,angle eccentricity=1.15,] {angle = m-2--ce--m0};
\coordinate (mm) at (\xs,-1.1);
\pic [angle radius=28mm,angle eccentricity=1.11,"$\varphi_2$"] {angle = mm--ce--m0};


%vec
\draw[->,red] (-1.2,2) ++ (-1,0) --++ (1,0) node[above,black,midway] {$\lambda$};
\draw[->,red] (-1.2,-2) ++ (-1,0) --++ (1,0);





    \end{tikzpicture}
\caption{Дифракция на решетке}
\label{fig:p193}
\end{figure}

На решетку~Р с периодом\footnote{Период решетки есть ширина щели плюс ширина непрозрачного промежутка.}~$d$ падает плоская волна красного света с длиной волны~$\lambda$. На экране~Э за решеткой наблюдается интерференционная картина~--- чередование максимумов и минимумов интерференции (места сильной освещенности отмечены красными полукругами~--- это центральный максимум~$O$ (в центре картины), а также первые и вторые максимумы, расположенные симметрично относительно максимума~$O$ и обозначенные цифрами~1 и~2). Направления, например, на первый и второй максимумы (относительно направления на центр картины) задают углы~$\varphi_1$ и~$\varphi_2$ соответственно.  

\textbf{Формула дифракционной решетки} описывает положения максимумов (с номерами~$m$) на экране:
\begin{equation}
    d\sin\varphi_m=m\lambda\quad(m=0,1,2,\ldots).
\end{equation}








 
\end{document}